% !TEX root = main.tex
% Theorem 4(i), fixed-perturbation guarantee: supplied routing under fixed off-direction query--key perturbations (K = 2, aligned),
% and the continuity proposition for small perturbations at every K. Input from appendix.tex after D.2.
% Numerical checks of every claim: Paper/research/planted_perturbation/theorem_checks.py.
\subsection{Supplied routing under query--key perturbations}
\label{app:perturbed-routing}

\Cref{thm:joint}(i) plants exact query and key directions. Here the gate head is perturbed off those directions,
\begin{equation}
 \mathbf W_Q=q\bm\Pi_Q+\bm\Delta_Q,\qquad \mathbf W_K=k\bm\Pi_K+\bm\Delta_K,\qquad \mathbf W_V=v\bm\Pi_V,
 \label{eq:planted-perturbation}
\end{equation}
with $\bm\Delta_Q,\bm\Delta_K$ fixed and $\langle\bm\Delta_K,\bm\Pi_K\rangle_F=0$ (a component along $\bm\Pi_K$ is a change of $k$).
Everything else is the seed of \cref{app:step-supplied-routing}, including the trained set $(\mathbf U,\mathbf W,q,k)$ and its metric.
We first compute the perturbed gate, then prove the fixed-perturbation guarantee in \cref{thm:joint}(i) for $K=2$ with an explicit constant radius, and finally
show that some radius exists at every $K$ and for both alignments.

\paragraph{The perturbed gate.} Without a value perturbation the readout features are those of \cref{lem:J-features};
only the gate changes.

\begin{lemma}[Perturbed gate]
\label{lem:perturbed-gate}
Under \cref{eq:planted-perturbation}, let $\bm\delta_s=\bm\Delta_Q\hat{\bm e}^{\mathrm c}_s$, $\beta_s=\langle\hat{\bm e}_{\mathrm h},\bm\delta_s\rangle$,
$\gamma_s=\langle\bm\delta_s,\bm\Delta_K\hat{\bm e}_\ell\rangle$, $\alpha_s=\langle\hat{\bm e}_{\mathrm h},\bm\Delta_K(\hat{\bm e}_0-\hat{\bm e}^{\mathrm c}_s)\rangle$,
$\omega_s=\langle\bm\delta_s,\bm\Delta_K(\hat{\bm e}_0-\hat{\bm e}^{\mathrm c}_s)\rangle$ and $\kappa=\initscale^2/\sqrt{d_{\mathrm h}}$. The prediction of $s_t$ with
$s_{t-1}=s$ and $b_t=b$ has logits $\initscale\mathbf U\bm e_s+\initscale v\,r^s_b\,\mathbf W\bm e_s$, where
\begin{equation}
 r^s_b=\sigma\bigl(-(-1)^ba_s-o_s\bigr),\qquad a_s=\kappa\bigl(k(q+\beta_s)+\gamma_s\bigr),\qquad o_s=\kappa(q\alpha_s+\omega_s),
 \label{eq:perturbed-gate}
\end{equation}
and
\begin{equation}
 r^s_0-r^s_1=\frac{\sinh(-a_s)}{\cosh a_s+\cosh o_s}.
 \label{eq:perturbed-contrast}
\end{equation}
With operator norms, $|\beta_s|\le\|\bm\Delta_Q\|$, $|\gamma_s|\le\|\bm\Delta_Q\|\|\bm\Delta_K\|$, $|\alpha_s|\le\sqrt2\|\bm\Delta_K\|$ and
$|\omega_s|\le\sqrt2\|\bm\Delta_Q\|\|\bm\Delta_K\|$.
\end{lemma}

\begin{proof}
As in \cref{lem:J-features}, block $1$ and the MLPs add zero. Since $\bm\Pi_Q\hat{\bm e}^{\mathrm c}_s=\hat{\bm e}_{\mathrm h}$,
$\bm\Pi_K\hat{\bm e}^{\mathrm c}_s=0$ and $\bm\Pi_K(\hat{\bm e}_0+(-1)^b\hat{\bm e}_\ell)=(-1)^b\hat{\bm e}_{\mathrm h}$, the query at $i_t$ and the
two keys are
\[
 \initscale(q\hat{\bm e}_{\mathrm h}+\bm\delta_s),\qquad \mathrm{key}_{i_t}=\initscale\bm\Delta_K\hat{\bm e}^{\mathrm c}_s,\qquad
 \mathrm{key}_{j_t}=\initscale\bigl((-1)^b(k\hat{\bm e}_{\mathrm h}+\bm\Delta_K\hat{\bm e}_\ell)+\bm\Delta_K\hat{\bm e}_0\bigr).
\]
The score difference $\mathrm{score}(j_t)-\mathrm{score}(i_t)$ is $\frac{\initscale^2}{\sqrt{d_{\mathrm h}}}$ times the inner product of the query with
$\mathrm{key}_{j_t}-\mathrm{key}_{i_t}$, which is $(-1)^ba_s+o_s$ because $\langle\hat{\bm e}_{\mathrm h},\bm\Delta_K\hat{\bm e}_\ell\rangle=\langle\bm\Delta_K,\bm\Pi_K\rangle_F=0$.
So $i_t$ gets weight $r^s_b$. The values are $\initscale v\hat{\bm e}^{\mathrm g}_s$ at $i_t$ and $0$ at $j_t$, so the residual and logits are those of
\cref{lem:J-features} with $r_b$ replaced by $r^s_b$. For \cref{eq:perturbed-contrast}, put $x=-a_s$, $o=o_s$; then
$\sigma(x-o)-\sigma(-x-o)=\frac{e^{x-o}-e^{-x-o}}{(1+e^{x-o})(1+e^{-x-o})}=\frac{2e^{-o}\sinh x}{2e^{-o}(\cosh x+\cosh o)}$. The bounds use
$\|\bm\delta_s\|\le\|\bm\Delta_Q\|$ and $\|\hat{\bm e}_0-\hat{\bm e}^{\mathrm c}_s\|=\sqrt2$.
\end{proof}

By \cref{eq:perturbed-contrast}, a state's letter contrast is positive exactly when $a_s<0$, whatever the offset; when $|o_s|\le|a_s|$ it is at
least $\frac12\tanh|a_s|$. The proof below therefore controls $a_s$ and $o_s$ state by state, relative to the planted strength
\[
 X:=\kappa\,q|k|.
\]

\paragraph{Setting of the fixed-perturbation guarantee.} $K=2$, $0<p<\frac12$, $h^\star=\log\frac{1-p}{p}$, an aligned start $q(0)k(0)<0$ with
$\frac{\initscale}{2}\le|q(0)|,|k(0)|,|v|\le2\initscale$, zero readouts, and
\begin{equation}
 \|\bm\Delta_Q\|\le\tfrac18|q(0)|,\qquad \|\bm\Delta_K\|\le\tfrac18|k(0)|,\qquad 0<\initscale\le p .
 \label{eq:perturbed-radius}
\end{equation}
Since $\|\bm\Pi_Q\|_F=\sqrt2$ and $\|\bm\Pi_K\|_F=1$, the Frobenius condition $\|\bm\Delta_X\|_F\le\rho\,|x_0|\,\|\bm\Pi_X\|_F$ implies
\cref{eq:perturbed-radius} for $\rho\le\frac{1}{8\sqrt2}$. Changing the signs of $\mathbf W_Q$ and $\mathbf W_K$ together, perturbations included, leaves every score
and \cref{eq:perturbed-radius} unchanged, so we take $q(0)>0>k(0)$; write $\delta=\|\bm\Delta_Q\|$ and $\eta_K=\|\bm\Delta_K\|$.

\begin{lemma}[Reduced flow for $K=2$]
\label{lem:perturbed-reduce}
Let $c_s=\initscale(U_{ss}-U_{s+1,s})$ and $g=2\initscale vw$. The step loss is
$L=\frac12\sum_{s\in\{0,1\}}L_s$ with $L_s=(1-p)\,\ell(m_{s0})+p\,\ell(m_{s1})$, $\ell(m)=\log(1+e^{-m})$, and margins
$m_{s0}=c_s+g\,r^s_0$, $m_{s1}=-c_s-g\,r^s_1$. With $E_{s0}=(1-p)\sigma(-m_{s0})$, $E_{s1}=p\,\sigma(-m_{s1})$,
$\Lambda_s=E_{s0}r^s_0(1-r^s_0)+E_{s1}r^s_1(1-r^s_1)$ and $\Gamma_s=E_{s0}r^s_0(1-r^s_0)-E_{s1}r^s_1(1-r^s_1)$, the flow is
\begin{gather*}
 \ddt{c_s}=\initscale^2(E_{s0}-E_{s1}),\qquad \ddt{g}=\frac{\initscale^2v^2}{2}\sum_s\bigl(E_{s0}r^s_0-E_{s1}r^s_1\bigr),\\
 \ddt{q}=-\frac{\kappa g}{2}\sum_s\bigl(\Lambda_sk+\Gamma_s\alpha_s\bigr),\qquad \ddt{k}=-\frac{\kappa g}{2}\sum_s\Lambda_s(q+\beta_s),
\end{gather*}
and the other two coordinates of $\mathbf U$ receive no gradient.
\end{lemma}

\begin{proof}
For $K=2$ the cross-entropy depends only on the difference of the two logits, which by \cref{lem:perturbed-gate} is $c_s+g\,r^s_b$ for the
prediction from state $s$; the target is $s+b$. The pair $(s_{t-1},b_t)$ is uniform $\times$ Bernoulli$(p)$ for every $t\ge2$, which gives $L$.
The coordinates $U_{ss}+U_{s+1,s}$ do not enter $L$. On $\mathbf U$ the Frobenius gradient gives $\ddt{c_s}=-2\initscale^2\pd{L}{c_s}$, and
$\|w(\mathbf I-\mathbf R)\|_F^2=4w^2$ gives $\ddt{g}=-\initscale^2v^2\pd{L}{g}$. Differentiating, $\pd{L_s}{c_s}=E_{s1}-E_{s0}$,
$\pd{L_s}{g}=E_{s1}r^s_1-E_{s0}r^s_0$, $\pd{L_s}{a_s}=g\Lambda_s$ and $\pd{L_s}{o_s}=g\Gamma_s$; with $\pd{a_s}{q}=\kappa k$, $\pd{a_s}{k}=\kappa(q+\beta_s)$,
$\pd{o_s}{q}=\kappa\alpha_s$ and $\pd{o_s}{k}=0$ this gives the gate equations.
\end{proof}

Note that $q^2-k^2$ is not conserved and $r^s_0+r^s_1\ne1$ when $o_s\ne0$; the next three lemmas replace the uses of these facts in
\cref{lem:J-align,lem:J-floor}.

\begin{lemma}[Statewise invariant region]
\label{lem:perturbed-invariant}
For all $t>0$: $q$ is nondecreasing and $k$ nonincreasing, so $q\ge q(0)$, $|k|\ge|k(0)|$ and $X$ is nondecreasing; $g>0$; and $c_s<h^\star$
for both $s$. Moreover, for both $s$,
\begin{equation}
 \begin{gathered}
 \tfrac{55}{64}X\le-a_s\le\tfrac{73}{64}X,\qquad |o_s|\le\tfrac{9\sqrt2}{64}X\le\tfrac15X,\\
 r^s_0\ge\tfrac12\ge r^s_1,\qquad r^s_0-r^s_1\ge\tfrac12\tanh|a_s|>0.
 \end{gathered}
 \label{eq:perturbed-gate-bounds}
\end{equation}
The flow exists for all $t\ge0$.
\end{lemma}

\begin{proof}
\emph{Gate bounds.} Wherever $q\ge q(0)\ge8\delta$ and $|k|\ge|k(0)|\ge8\eta_K$, \cref{lem:perturbed-gate} gives
$-a_s=\kappa(|k|(q+\beta_s)-\gamma_s)\in\kappa q|k|\,[1-\frac18-\frac1{64},\,1+\frac18+\frac1{64}]$ and
$|o_s|\le\kappa\sqrt2\eta_K(q+\delta)\le\kappa\sqrt2\frac{|k|}8\cdot\frac{9q}8$, which is \cref{eq:perturbed-gate-bounds} for $a_s,o_s$. Then $|o_s|<|a_s|$,
so $r^s_0=\sigma(|a_s|-o_s)\ge\frac12\ge\sigma(-|a_s|-o_s)=r^s_1$, and \cref{eq:perturbed-contrast} with $\cosh o_s\le\cosh a_s$ gives the contrast bound.

\emph{Gate monotonicity on intervals with $g\ge0$.} Since $|\Gamma_s|\le\Lambda_s$, $|\alpha_s|\le\sqrt2\eta_K$ and $k<0$,
\[
 \begin{aligned}
 \ddt{q}&=\frac{\kappa g}{2}\sum_s\bigl(\Lambda_s|k|-\Gamma_s\alpha_s\bigr)
 \ge\frac{\kappa g}{2}\bigl(|k|-\sqrt2\eta_K\bigr)\sum_s\Lambda_s,\\
 \ddt{|k|}&=\frac{\kappa g}{2}\sum_s\Lambda_s(q+\beta_s)
 \ge\frac{\kappa g}{2}(q-\delta)\sum_s\Lambda_s.
 \end{aligned}
\]
Both are nonnegative while $|k|>\sqrt2\eta_K$ and $q>\delta$. These strict conditions hold at $t=0$, and by the displayed signs they cannot first fail
inside an interval on which $g\ge0$; so on any such interval $q\ge q(0)$, $|k|\ge|k(0)|$, both are monotone, and the gate bounds hold.

\emph{Readout region.} At $t=0$, $g=0$, $c_s=0$ and $E_{s0}=\frac{1-p}2>\frac p2=E_{s1}$, so $\ddt{g}(0)>0$ by $r^s_0>r^s_1$; hence
$t^\ast=\inf\{t>0:g(t)\le0\text{ or }\max_sc_s(t)\ge h^\star\}>0$. Suppose $t^\ast<\infty$; on $[0,t^\ast]$, $g\ge0$ and the gate bounds hold.
If $g(t^\ast)=0$, then $\ddt{g}(t^\ast)\le0$. But at $g=0$, $E_{s0}r^s_0-E_{s1}r^s_1=(1-p)\sigma(-c_s)r^s_0-p\,\sigma(c_s)r^s_1>0$, because
$p\,\sigma(c_s)\le(1-p)\sigma(-c_s)$ for $c_s\le h^\star$ and $r^s_0>r^s_1$; a contradiction. Otherwise $g(t^\ast)>0$ and $c_s(t^\ast)=h^\star$ for some
$s$, so $\ddt{c_s}(t^\ast)\ge0$. But $\sigma(-h^\star)=p$, so $E_{s0}=(1-p)\sigma(-h^\star-gr^s_0)<p(1-p)\le p\,\sigma(h^\star+gr^s_1)=E_{s1}$, a contradiction.
So $t^\ast=\infty$.

\emph{Existence.} $|\ddt{c_s}|\le\initscale^2$ and $|\ddt{g}|\le\initscale^2v^2$, so $c_s,g$ grow at most linearly; given $g$, the $(q,k)$ equations are
linear-growth, so Gr\"onwall excludes blow-up in finite time.
\end{proof}

\begin{lemma}[Balance and readout floor]
\label{lem:perturbed-floor}
Let $Q=g-\frac{v^2}4(c_0+c_1)$, $\mu=\frac{p(1-2p)}3$, $g_{\mathrm{low}}=\frac{v^2h^\star}{16}$ and $T_h=\frac{h^\star}{2\initscale^2\mu}$.
On every interval $[0,t]$ on which $X\le1$, $Q$ is nondecreasing and $Q\ge0$, and $g(t)\ge g_{\mathrm{low}}$ if $t\ge T_h$.
\end{lemma}

\begin{proof}
\emph{Balance.} Write $E_{s0}r^s_0-E_{s1}r^s_1=\bar r_s(E_{s0}-E_{s1})+\frac12(r^s_0-r^s_1)(E_{s0}+E_{s1})$ with $\bar r_s=\frac12(r^s_0+r^s_1)$.
Since $|\partial_o(\sigma(-a-o)+\sigma(a-o))|\le\frac12$, $|\bar r_s-\frac12|\le\frac{|o_s|}4$, and so
\[
 \begin{aligned}
 \ddt{Q}&=\frac{\initscale^2v^2}2\sum_s\Bigl[
 \bigl(\bar r_s-\tfrac12\bigr)(E_{s0}-E_{s1})
 +\tfrac12(r^s_0-r^s_1)(E_{s0}+E_{s1})\Bigr]\\
 &\ge\frac{\initscale^2v^2}4\sum_s(E_{s0}+E_{s1})
 \Bigl(r^s_0-r^s_1-\frac{|o_s|}2\Bigr).
 \end{aligned}
\]
For $X\le1$, \cref{eq:perturbed-gate-bounds} gives $|a_s|\le\frac{73}{64}$, so $r^s_0-r^s_1\ge\frac12\tanh|a_s|\ge0.35\,|a_s|\ge0.3X\ge\frac{|o_s|}2$.
Hence $\ddt{Q}\ge0$, and $Q(0)=0$.

\emph{Floor.} For $c\le h^\star$, $0\le g\le1$ and $r_0,r_1\in[0,1]$, $E_0-E_1\ge f(c+g)$ with $f(y)=(1-p)\sigma(-y)-p\,\sigma(y)$, which is
decreasing, has $|\dd{f}{y}|\le\frac14$, $f(h^\star)=0$, and
$f(\frac{h^\star}2)=\frac{\sqrt{p(1-p)}(1-2p)}{1+2\sqrt{p(1-p)}}\ge\mu$. So $\ddt{c_s}\ge\initscale^2\mu$ when $c_s+g\le\frac{h^\star}2$, and
$\ddt{c_s}\ge-\frac{\initscale^2g}4$ always. By \cref{eq:perturbed-radius} and $h^\star\le\frac{1-2p}{p}$,
$g_{\mathrm{low}}\le\frac{\initscale^2h^\star}4\le\min\{1,\frac{3h^\star}8,2\mu\}$.

Suppose $t\ge T_h$, $X\le1$ on $[0,t]$ and $g(t)<g_{\mathrm{low}}$. Let $s_0=\sup\{s'\le t:g(s')\ge g_{\mathrm{low}}\}$, or $s_0=0$ if there is none.
On $(s_0,t]$, $g<g_{\mathrm{low}}$ and $Q\ge0$ give $c_0+c_1\le\frac{4g}{v^2}<\frac{h^\star}4$. Some state then has $c_s<\frac{h^\star}8$ and
$c_s+g<\frac{h^\star}2$, so with the other state's lower bound
\[
 \ddt{}(c_0+c_1)\ge\initscale^2\bigl(\mu-\tfrac{g_{\mathrm{low}}}4\bigr)\ge\tfrac12\initscale^2\mu,\qquad
 \ddt{g}=\ddt{Q}+\tfrac{v^2}4\ddt{}(c_0+c_1)>0 .
\]
If $s_0>0$, then $g(s_0)=g_{\mathrm{low}}$ and $g$ increases on $(s_0,t]$, contradicting $g(t)<g_{\mathrm{low}}$. If $s_0=0$, then
$c_0+c_1\ge\frac12\initscale^2\mu\,T_h=\frac{h^\star}4$ at $t$, a contradiction.
\end{proof}

\begin{lemma}[Error floor and gate growth]
\label{lem:perturbed-grow}
Let $z=\frac p4\log2$ and $e_\ast=\frac{p^2}{12}$. If $L>z$, then $\sum_{s,b}E_{sb}\ge e_\ast$. If moreover $t\ge T_h$ and $X\le1$, then
$\ddt{X}\ge\frac{\kappa\,g_{\mathrm{low}}\,p^2}{105}\,X$.
\end{lemma}

\begin{proof}
Each of the four (state, letter) cases has weight $w_{sb}=\frac{\pi_b}2\in[\frac p2,\frac12]$. If $L>z$, some case has $w_{sb}\ell(m_{sb})\ge\frac L4$,
so $\ell(m_{sb})>\frac z2$; since $\sigma(-m)=1-e^{-\ell(m)}$ and $\frac z2<0.05$, $E_{sb}\ge p(1-e^{-z/2})\ge0.97\,p\frac z2\ge e_\ast$.

For the growth, \cref{lem:perturbed-invariant} gives $|k|-\sqrt2\eta_K\ge(1-\frac{\sqrt2}8)|k|$ and $q-\delta\ge\frac78q$, so
\[
 \ddt{(q|k|)}=|k|\ddt{q}+q\ddt{|k|}\ge\frac{\kappa g}2\cdot\frac45(q^2+k^2)\sum_s\Lambda_s\ge\frac45\kappa g\,q|k|\sum_s\Lambda_s .
\]
For $X\le1$ every attention argument satisfies $|a_s|+|o_s|\le\frac{73+9\sqrt2}{64}<1.34$, so $r(1-r)\ge\dd{\sigma(x)}{x}\big|_{x=1.34}\ge\frac17$ and
$\sum_s\Lambda_s\ge\frac17\sum_{s,b}E_{sb}\ge\frac{e_\ast}7$. With $g\ge g_{\mathrm{low}}$ from \cref{lem:perturbed-floor},
$\ddt{X}\ge\frac45\cdot\frac17\cdot\frac{p^2}{12}\kappa g_{\mathrm{low}}X$.
\end{proof}

\begin{lemma}[Statewise descent]
\label{lem:perturbed-descent}
If $X\ge1$, then $\ddt{L}\le-\frac{\initscale^4}{6732}\bigl(\sum_{s,b}E_{sb}\bigr)^2$.
\end{lemma}

\begin{proof}
The readout part of the flow alone dissipates $2\initscale^2\sum_s(\pd{L}{c_s})^2+\initscale^2v^2(\pd{L}{g})^2
=\frac{\initscale^2}2D^2+\frac{\initscale^2v^2}4G^2$, with $D^2=\sum_s(E_{s0}-E_{s1})^2$ and $G=\sum_s(E_{s0}r^s_0-E_{s1}r^s_1)$; the gate part only adds
dissipation. For $X\ge1$, \cref{eq:perturbed-gate-bounds} gives $r^s_0-r^s_1\ge b:=\frac12\tanh\frac{55}{64}\ge\frac13$. Write $S_0=\sum_sE_{s0}$.
Since $E_{s0}r^s_0-E_{s1}r^s_1=E_{s0}(r^s_0-r^s_1)+r^s_1(E_{s0}-E_{s1})$, $G\ge bS_0-\sqrt2D$. If $D\ge\frac{bS_0}{2\sqrt2}$ the first term is at least
$\frac{\initscale^2b^2S_0^2}{16}$; otherwise $G\ge\frac{bS_0}2$ and the second is at least $\frac{\initscale^2v^2b^2S_0^2}{16}$. As $v\le1$, the dissipation is at
least $\frac{\initscale^2v^2b^2}{16}\max\{S_0,D\}^2$, and $\sum_{s,b}E_{sb}\le2S_0+\sqrt2D\le(2+\sqrt2)\max\{S_0,D\}$. Use $v^2\ge\frac{\initscale^2}4$ and
$b^2\ge\frac19$.
\end{proof}

\begin{proof}[Fixed-perturbation guarantee of \cref{thm:joint}(i)]
By \cref{lem:perturbed-invariant} the flow exists, stays in the aligned region and $X$ is nondecreasing, with
$X(0)\ge\frac{\initscale^4}{4\sqrt{d_{\mathrm h}}}$ and $X(0)\le\frac{4\initscale^4}{\sqrt{d_{\mathrm h}}}<1$. Let $t_3=\inf\{t:L(t)\le z\}$.

\emph{Growth.} Let $t_1=\inf\{t:X(t)\ge1\}$. On $[T_h,\min\{t_1,t_3\})$, \cref{lem:perturbed-grow} gives exponential growth of $X$ at rate
$\frac{\kappa g_{\mathrm{low}}p^2}{105}\ge\frac{\initscale^4h^\star p^2}{6720\sqrt{d_{\mathrm h}}}$, so
$\min\{t_1,t_3\}\le T_h+\frac{6720\sqrt{d_{\mathrm h}}}{\initscale^4h^\star p^2}\log\frac{4\sqrt{d_{\mathrm h}}}{\initscale^4}$.

\emph{Descent.} If $t_1<t_3$, then $X\ge1$ on $[t_1,\infty)$, and \cref{lem:perturbed-descent,lem:perturbed-grow} give
$\ddt{L}\le-\frac{\initscale^4e_\ast^2}{6732}$ while $L>z$; since $L(t_1)\le\log2$, $t_3-t_1\le\frac{6732\cdot144\log2}{\initscale^4p^4}$.

\emph{Margins.} If $L\le z$, every case has $\ell(m_{sb})\le\frac{z}{w_{sb}}\le\frac{\log2}2<\log2$, so every margin is positive; $L$ never increases,
so this persists for all $t\ge t_3$. After $t_3$, all margins are positive, so
$\ell(m)\le2\sigma(-m)$ and $L\le\sum_{s,b}E_{sb}$.
If $t_1<\infty$, put $t_4=\max\{t_1,t_3\}$. For $t\ge t_4$,
\cref{lem:perturbed-descent} applies and gives
\[
 \ddt{L}\le-\frac{\initscale^4}{6732}L^2,\qquad
 L(t)\le\frac{6732}{\initscale^4(t-t_4)}\longrightarrow0
 \quad(t>t_4).
\]
Otherwise $X<1$ for all time. If the nonincreasing loss had limit
$L_\infty>0$, then for $t\ge\max\{T_h,t_3\}$,
\cref{lem:perturbed-floor} and the growth calculation in
\cref{lem:perturbed-grow} would give
\[
 \sum_s\Lambda_s\ge\frac17\sum_{s,b}E_{sb}\ge\frac{L_\infty}{7},\qquad
 \ddt{X}\ge\frac{4\kappa g_{\mathrm{low}}L_\infty}{35}X.
\]
Exponential growth contradicts $X<1$. Hence $L_\infty=0$ in this case too.

\emph{Time.} $T_h=\frac{3h^\star}{2\initscale^2p(1-2p)}\le\frac{3}{2\initscale^2p^2}$, so
\[
 t_3\le\frac{3}{2\initscale^2p^2}+\frac{6720\sqrt{d_{\mathrm h}}}{\initscale^4h^\star p^2}\log\frac{4\sqrt{d_{\mathrm h}}}{\initscale^4}+\frac{7\cdot10^5}{\initscale^4p^4}
 \le\frac{C}{\initscale^4}\Bigl(\frac{\sqrt{d_{\mathrm h}}\log\frac{d_{\mathrm h}}{\initscale}}{p^2\min\{1,h^\star\}}+\frac1{p^4}\Bigr).\qedhere
\]
\end{proof}

The radius in \cref{eq:perturbed-radius} is a constant: it does not shrink with $\initscale$, $d_{\mathrm h}$ or $n$. The fixed perturbation is also
relatively suppressed during training, since $q$ and $|k|$ only grow.

\paragraph{Every $K$, both alignments: some radius.} Without an explicit radius, continuity extends \cref{thm:joint}(i) to small fixed perturbations
of the whole gate head.

\begin{proposition}[Small fixed perturbations]
\label{prop:perturbed-continuity}
Take any seed and $\initscale$ for which \cref{thm:joint}(i) applies, aligned or misaligned, and any fixed matrices $\bm\Delta_Q,\bm\Delta_K,\bm\Delta_V$ added to
$\mathbf W_Q,\mathbf W_K,\mathbf W_V$, with the same trained set and metric. There are $\rho_\star>0$ and $T_\star<\infty$, depending on the seed, such that
$\|\bm\Delta_Q\|_F+\|\bm\Delta_K\|_F+\|\bm\Delta_V\|_F\le\rho_\star$ implies $\Lstep\le\frac{p_{\min}\log2}{2K}$ by $T_\star$ and positive margins on every
transition for all $t\ge T_\star$.
\end{proposition}

\begin{proof}
Let $\eta_\star=\frac{p_{\min}\log2}{2K}$. By \cref{thm:joint}(i) the unperturbed loss tends to $0$, so $\Lstep\le\eta_\star/2$ at some $T_\star$. The
population loss is real-analytic in the trained parameters and in $\bm\Delta=(\bm\Delta_Q,\bm\Delta_K,\bm\Delta_V)$, so the flow depends continuously on
$\bm\Delta$ on the compact interval $[0,T_\star]$, and so does the loss at $T_\star$; this gives $\rho_\star$. Each (state, letter) case has weight at
least $\frac{p_{\min}}K$, so $\Lstep\le\eta_\star$ makes each cross-entropy at most $\frac{\log2}2$ and each correct-class probability at least
$2^{-1/2}>\frac12$. Gradient flow does not increase $\Lstep$, so the margins persist.
\end{proof}

Unlike \cref{thm:joint}(i), \cref{prop:perturbed-continuity} gives no size for $\rho_\star$; a Gr\"onwall bound for it can be exponentially small
in the $\initscale^{-4}$ horizon.

\paragraph{Completion under a gate floor, for any query--key dynamics.} The remaining gap for unplanted query--key maps is the gate, not
the readout: once every state's letter contrast is bounded below, the readouts finish in polynomial time while the query--key maps keep training.

\begin{lemma}[Completion under a statewise gate floor]
\label{lem:gate-floor-completion}
Take the supplied two-site mask, embeddings and value map $\mathbf W_V=v\bm\Pi_V$, $v\ne0$, of \cref{app:step-supplied-routing}, with any $K\ge2$ and
$p\in(0,1)$, zero initial readouts, and arbitrary query and key maps (and relative biases) $\bm\theta_{\mathrm{QK}}$. Train $(\mathbf U,\mathbf W,\bm\theta_{\mathrm{QK}})$
jointly by gradient flow with the Frobenius metric on $\mathbf U$ and $\mathbf W\in\R(\mathbf I-\mathbf R)$ and any block-separate metric on
$\bm\theta_{\mathrm{QK}}$. Let $M=\frac{K\log K}{p_{\min}}$, $\eta_\star=\frac{p_{\min}\log2}{2K}$ and
\[
 T_{\mathrm{fin}}=\frac{8K(1+v^2)(1+M)^2}{\initscale^2v^2\tau_0^2\,\eta_\star}.
\]
If $\min_s(r^s_0-r^s_1)\ge\tau_0>0$ on $[T_0,t]$ for every $t<T_0+T_{\mathrm{fin}}$ with $\Lstep(t)>\eta_\star$, then $\Lstep\le\eta_\star$ by $T_0+T_{\mathrm{fin}}$,
and every transition has positive margin at all later times, whatever the gate does afterwards. The same holds with $r^s_1-r^s_0$ in place of
$r^s_0-r^s_1$.
\end{lemma}

\begin{proof}
The value at the letter site is $\mathbf W_V\initscale(\hat{\bm e}_0\pm\hat{\bm e}_\ell)=0$, so for any scores the prediction from $s_{t-1}=s$ with letter $b$ has logits
$\bm z_{sb}=\initscale\mathbf U\bm e_s+\initscale v\,r^s_b\,\mathbf W\bm e_s$, linear in $(\mathbf U,\mathbf W)$ at fixed attention.
At a time $t$ with the floor, put $m_s=\frac12(r^s_0+r^s_1)$, $w_\star=\frac{2}{\initscale v\tau_0}$, $\mathbf W_\star=w_\star(\mathbf I-\mathbf R)$ and
$\mathbf U_\star\bm e_s=-v\,w_\star m_s(\bm e_s-\bm e_{s+1})$. These produce logits
$\bm z^\star_{sb}=\frac{2}{\tau_0}(r^s_b-m_s)(\bm e_s-\bm e_{s+1})=\pm\frac{r^s_0-r^s_1}{\tau_0}(\bm e_s-\bm e_{s+1})$, with sign $+$ for $b=0$ (target $s$) and $-$
for $b=1$ (target $s+1$). So the target logit exceeds every other by at least $1$. Since $m_s\le1$,
$\|\mathbf U_\star\|_F^2+\|\mathbf W_\star\|_F^2\le2K(1+v^2)w_\star^2=:C_\star^2$.

Let $\pi_{sb}\ge\frac{p_{\min}}K$ be the case weights and $P_{sb}$ the probability of the target. By linearity, the derivative of $\Lstep$ along
$(\mathbf U_\star,\mathbf W_\star)$ is $\sum\pi_{sb}\langle\softmax(\bm z_{sb})-\bm e_{y},\bm z^\star_{sb}\rangle\le-\sum\pi_{sb}(1-P_{sb})$, so by Cauchy--Schwarz
$\|\nabla_{\mathbf U,\mathbf W}\Lstep\|\ge\frac1{C_\star}\sum\pi_{sb}(1-P_{sb})$. Since $\Lstep\le\Lstep(0)=\log K$, every case has $-\log P_{sb}\le M$, and because
$-\log P/(1-P)$ decreases in $P$, $-\log P_{sb}\le\frac{M}{1-e^{-M}}(1-P_{sb})\le(1+M)(1-P_{sb})$. Hence $\Lstep\le(1+M)C_\star\|\nabla_{\mathbf U,\mathbf W}\Lstep\|$ and,
as the query--key block only adds dissipation,
\[
 \ddt{\Lstep}\le-\|\nabla_{\mathbf U,\mathbf W}\Lstep\|^2\le-\frac{\Lstep^2}{(1+M)^2C_\star^2},\qquad
 \Lstep(t)\le\frac{(1+M)^2C_\star^2}{t-T_0}.
\]
This reaches $\eta_\star$ by $T_0+T_{\mathrm{fin}}$. Then each cross-entropy is at most $\frac{\log2}2$, and since $\Lstep$ never increases, every margin stays
positive. For the reversed orientation use $-w_\star$.
\end{proof}

With $v=\Theta(\initscale)$, $T_{\mathrm{fin}}=O\bigl(K^4\log^2K/(p_{\min}^3\initscale^4\tau_0^2)\bigr)$. The lemma reduces convergence from unplanted query--key
maps to one statement: gradient flow must create a statewise gate floor $\tau_0\ge1/\mathrm{poly}$ in polynomial time and keep it until
completion. That statement is not proved here. \Cref{app:planted-perturbation} measures it.

\paragraph{Toward a random query--key bootstrap.} The next three results describe the start of that missing step. Keep the
supplied mask, embeddings, $\mathbf W_V=v\bm\Pi_V$ and readouts, and let $\mathbf W_Q,\mathbf W_K$ be arbitrary full matrices. For any scores write
$r^s_b=\sigma(-(-1)^ba_s-o_s)$ as in \cref{lem:perturbed-gate}, $d_{b,s}=r^s_b(1-r^s_b)$, $P_{sb}=\softmax(\bm z_{sb})$ and
\[
 \zeta_{0,s}=1-P_{s0;s}+P_{s0;s+1},\qquad \zeta_{1,s}=1+P_{s1;s}-P_{s1;s+1},
\]
both strictly positive for finite logits.

\begin{lemma}[A common useful-score direction]
\label{lem:score-direction}
For every $K\ge2$, every readout $\mathbf U$ and all scores,
\[
 \pd{\Lstep}{a_s}=\frac{g}{2K}\bigl[(1-p)\zeta_{0,s}d_{0,s}+p\,\zeta_{1,s}d_{1,s}\bigr],\qquad
 \pd{\Lstep}{o_s}=\frac{g}{2K}\bigl[(1-p)\zeta_{0,s}d_{0,s}-p\,\zeta_{1,s}d_{1,s}\bigr],
\]
so every $\pd{\Lstep}{a_s}$ has the sign of $g$. At zero readouts, with $\bar r_b=\frac1K\sum_sr^s_b$,
$\ddt{g}(0)=\frac{\initscale^2v^2}{K}\bigl[(1-p)\bar r_0-p\,\bar r_1\bigr]$; if $p<\frac12$ and $|a_s(0)|+|o_s(0)|\le\delta<h^\star$ for every $s$, then
$\ddt{g}(0)\ge\frac{\initscale^2v^2}{K}(\sigma(-\delta)-p)>0$.
\end{lemma}

\begin{proof}
The logits are $\bm z_{sb}=\initscale\mathbf U\bm e_s+\frac g2r^s_b\bm x_s$ with $\bm x_s=\bm e_s-\bm e_{s+1}$, so
$\partial_{r^s_b}\mathrm{CE}=\frac g2\langle P_{sb}-\bm e_{s+b},\bm x_s\rangle=\frac g2\bigl(P_{sb;s}-P_{sb;s+1}-(1-2b)\bigr)$, which is $-\frac g2\zeta_{0,s}$ for
$b=0$ and $\frac g2\zeta_{1,s}$ for $b=1$. With $\partial_{a_s}r^s_0=-d_{0,s}$, $\partial_{a_s}r^s_1=d_{1,s}$, $\partial_{o_s}r^s_b=-d_{b,s}$ and case weights
$\frac{\pi_b}K$ this gives both derivatives. At zero readouts $P_{sb}$ is uniform, so $\partial_w\mathrm{CE}=\mp\initscale v\,r^s_b$ for $b=0,1$, and
$\ddt{w}=-\frac1{2K}\partial_w\Lstep$ on the line $\R(\mathbf I-\mathbf R)$ gives $\ddt{g}(0)$. Finally $r^s_0\ge\sigma(-\delta)$ and $r^s_1\le1-\sigma(-\delta)$, and
$\sigma(-\delta)>p$ exactly when $\delta<h^\star$.
\end{proof}

Write $\kappa=\initscale^2/\sqrt{d_{\mathrm h}}$ and let $\bm q_s=\mathbf W_Q\hat{\bm e}^{\mathrm c}_s$, $\bm h=\mathbf W_K\hat{\bm e}_\ell$, $\bm k_0=\mathbf W_K\hat{\bm e}_0$ and
$\bm k_s=\mathbf W_K\hat{\bm e}^{\mathrm c}_s$ be the columns the head reads. As in the proof of \cref{lem:perturbed-gate},
\begin{equation}
 a_s=\kappa\langle\bm q_s,\bm h\rangle,\qquad o_s=\kappa\langle\bm q_s,\bm k_0-\bm k_s\rangle .
 \label{eq:factor-scores}
\end{equation}

\begin{lemma}[Factorized query--key flow]
\label{lem:factor-flow}
Under Frobenius gradient flow of $\mathbf W_Q,\mathbf W_K$, with $A_s=\pd{\Lstep}{a_s}$ and $B_s=\pd{\Lstep}{o_s}$,
\begin{gather*}
 \ddt{\bm q_s}=-\kappa\bigl[A_s\bm h+B_s(\bm k_0-\bm k_s)\bigr],\qquad
 \ddt{\bm h}=-\kappa\sum_sA_s\bm q_s,\\
 \ddt{\bm k_0}=-\kappa\sum_sB_s\bm q_s,\qquad
 \ddt{\bm k_s}=\kappa B_s\bm q_s,\\
 \ddt{a_s}=-\kappa^2\Bigl[A_s\|\bm h\|^2+B_s\langle\bm k_0-\bm k_s,\bm h\rangle+\sum_jA_j\langle\bm q_s,\bm q_j\rangle\Bigr],
\end{gather*}
and the balance $\sum_s\|\bm q_s\|^2-\bigl(\|\bm h\|^2+\|\bm k_0\|^2+\sum_s\|\bm k_s\|^2\bigr)$ is conserved; it replaces $q^2-k^2$.
\end{lemma}

\begin{proof}
The head reads $\mathbf W_Q,\mathbf W_K$ only through these columns, which sit on orthonormal input directions, so Frobenius flow is Euclidean flow on
them, and the chain rule through \cref{eq:factor-scores} gives the column equations and $\ddt{a_s}$. Both
$\ddt{}\sum_s\|\bm q_s\|^2$ and $\ddt{}\bigl(\|\bm h\|^2+\|\bm k_0\|^2+\sum_s\|\bm k_s\|^2\bigr)$ equal $-2\sum_s(A_sa_s+B_so_s)$.
\end{proof}

By \cref{lem:score-direction} the direct score gradient points every $a_s$ towards a useful gate as soon as $g>0$. In parameter space,
the Gram couplings $\langle\bm q_s,\bm q_j\rangle$ and the offset term can oppose it. At initialization they are small for Gaussian columns
of large head dimension.

\begin{proposition}[Initial gate acceleration from Gaussian query--key maps]
\label{prop:gaussian-acceleration}
Let $p<\frac12$, zero readouts, and $\bm q_s(0),\bm h(0),\bm k_0(0),\bm k_s(0)$ independent $\mathcal N(0,\sigma^2\mathbf I_{d_{\mathrm h}})$. Let
$0<\delta<1$, $u=\log\frac{16K^2}\delta$, $d_{\mathrm h}\ge2048K^2u$ and
$\initscale^2\sigma^2\bigl(1+\sqrt{2u/d_{\mathrm h}}\bigr)(\sqrt{2u}+2\sqrt u)<\min\{1,h^\star\}$. With probability at least $1-\delta$,
$\ddt{g}(0)>0$ and, for every state, $\ddt{a_s}(0)=0$ and $\frac{\mathrm d^2a_s}{\mathrm dt^2}(0)<0$.
\end{proposition}

\begin{proof}
Let $A=\sqrt{d_{\mathrm h}}+\sqrt{2u}\le2\sqrt{d_{\mathrm h}}$. Gaussian norm concentration gives $\|\bm q_s\|,\|\bm h\|\le\sigma A$, and the $\chi^2$ lower
tail gives $\|\bm h\|^2\ge\frac{d_{\mathrm h}\sigma^2}4$ (as $u\le\frac{9d_{\mathrm h}}{64}$). Conditioning on one factor of each inner product,
$|\langle\bm q_s,\bm h\rangle|,|\langle\bm q_s,\bm q_j\rangle|\le\sigma^2A\sqrt{2u}\le3\sigma^2\sqrt{d_{\mathrm h}u}$ ($j\ne s$) and
$|\langle\bm q_s,\bm k_0-\bm k_s\rangle|,|\langle\bm h,\bm k_0-\bm k_s\rangle|\le2\sigma^2A\sqrt u\le4\sigma^2\sqrt{d_{\mathrm h}u}$. A union bound
over at most $(K^2+6K+2)$ events of probability $e^{-u}$ each leaves probability at least $1-\delta$. On this event
$|a_s(0)|+|o_s(0)|\le\kappa\sigma^2A(\sqrt{2u}+2\sqrt u)<\min\{1,h^\star\}$ by \cref{eq:factor-scores}, so \cref{lem:score-direction} gives $\ddt{g}(0)>0$. At $t=0$, $g=0$, so $A_s=B_s=0$,
$\ddt{a_s}(0)=0$, $\ddt{A_s}(0)=\frac{\ddt{g}(0)}{2K}[(1-p)d_{0,s}+p\,d_{1,s}]>0$ and $|\ddt{B_s}(0)|\le\ddt{A_s}(0)$. Since every score argument
has magnitude below $1$, $d_{b,s}\in(\frac18,\frac14]$ and $\ddt{A_j}(0)\le2\ddt{A_s}(0)$. Differentiating \cref{lem:factor-flow} at $t=0$,
\[
 \begin{aligned}
 -\frac{1}{\kappa^2}\frac{\mathrm d^2a_s}{\mathrm dt^2}(0)
 &=\ddt{A_s}\bigl(\|\bm h\|^2+\|\bm q_s\|^2\bigr)
 +\sum_{j\ne s}\ddt{A_j}\langle\bm q_s,\bm q_j\rangle\\
 &\quad+\ddt{B_s}\langle\bm k_0-\bm k_s,\bm h\rangle\\
 &\ge\ddt{A_s}\,\sigma^2\Bigl(\frac{d_{\mathrm h}}4-(6K-2)\sqrt{d_{\mathrm h}u}\Bigr)>0.
 \end{aligned}
\]
since $d_{\mathrm h}>16(6K-2)^2u$.
\end{proof}

The next lemma makes the missing step quantitative. It reaches the gate floor in polynomial time under explicit trajectory conditions,
without any proximity to $\bm\Pi_Q,\bm\Pi_K$.

\begin{lemma}[Conditional gate hitting time]
\label{lem:gate-hitting}
Let $0<a$, and suppose that on an interval $[t_0,t_1]$ of the flow of \cref{lem:factor-flow}: $0<g_0\le g\le G$; $|o_s|\le B$ for all $s$;
and $\|\bm h\|^2\ge h_0>0$. Let $A_0=\max\{0,\max_sa_s(t_0)\}$, $d_\ast=\sigma'(B+\max\{A_0,a\})$, $\psi=(1+e^{Ga/4})^{-1}$ and
$R_\ast=(2p_{\min}\psi d_\ast)^{-1}$. Suppose further that at every time and state with $a_s\ge-a$,
\[
 \sum_{j\ne s}|\langle\bm q_s,\bm q_j\rangle|\le\frac{h_0}{4R_\ast},\qquad |\langle\bm k_0-\bm k_s,\bm h\rangle|\le\frac{h_0}4 .
\]
Then every $a_s$ is at most $-a$ from time
\[
 T_{\mathrm{gate}}=t_0+\frac{4Kd_{\mathrm h}(A_0+a)}{\initscale^4h_0g_0p_{\min}\psi d_\ast}
\]
on, as long as $T_{\mathrm{gate}}\le t_1$, and remains so until $t_1$.
\end{lemma}

\begin{proof}
Fix a time and a state with $a_s\ge-a$. The adjacent-class margins of its two transitions are $m_{0,s}=z_{s0;s}-z_{s0;s+1}$ and
$m_{1,s}=z_{s1;s+1}-z_{s1;s}$; the readout $\mathbf U$ contributes equally to both, so $m_{0,s}+m_{1,s}=g(r^s_0-r^s_1)$. If $a_s\ge0$ the contrast is
nonpositive; if $-a\le a_s<0$ it is at most $\frac{|a_s|}2\le\frac a2$ because $\sigma$ is $\frac14$-Lipschitz. So one of the two margins is at most
$\frac{Ga}4$, and that transition's target probability is at most $\sigma(\frac{Ga}4)$, giving $\zeta_{b,s}\ge1-P_{sb;s+b}\ge\psi$. As long as
$-a\le a_s\le A_0$, the score arguments $-o_s\mp a_s$ have magnitude at most $B+\max\{A_0,a\}$, so $d_{b,s}\ge d_\ast$; the bound below then
makes $a_s$ decrease, so by continuity $a_s$ never exceeds $A_0\ge a_s(t_0)$ and the bound applies whenever $a_s\ge-a$. By \cref{lem:score-direction},
$A_s\ge\frac{g}{2K}p_{\min}\psi d_\ast$, while for every $j$, $A_j\le\frac{g}{2K}\cdot2\cdot\frac14=\frac g{4K}\le R_\ast A_s$ and $|B_s|\le A_s$. In
\cref{lem:factor-flow} the diagonal term $A_s\|\bm q_s\|^2$ is nonnegative, so the bracket is at least
$A_sh_0-R_\ast A_s\frac{h_0}{4R_\ast}-A_s\frac{h_0}4=\frac{A_sh_0}2$, and
\[
 \ddt{a_s}\le-\frac{\initscale^4}{d_{\mathrm h}}\cdot\frac{h_0}2\cdot\frac{g_0p_{\min}\psi d_\ast}{2K}.
\]
Each $a_s$ starts at most $A_0$ and must fall by at most $A_0+a$, which gives $T_{\mathrm{gate}}$; at $a_s=-a$ the derivative is negative, so
$a_s$ does not return above $-a$ before $t_1$.
\end{proof}

\begin{corollary}[Conditional polynomial bootstrap]
\label{cor:conditional-bootstrap}
If the hypotheses of \cref{lem:gate-hitting} hold on $[t_0,T_{\mathrm{gate}}+T_{\mathrm{fin}}]$, with $T_{\mathrm{fin}}$ from
\cref{lem:gate-floor-completion} at $\tau_0=\frac{\sinh a}{\cosh a+\cosh B}$, then every transition has positive margin from
$T_{\mathrm{gate}}+T_{\mathrm{fin}}$ on.
\end{corollary}

\begin{proof}
On $[T_{\mathrm{gate}},T_{\mathrm{gate}}+T_{\mathrm{fin}}]$, $a_s\le-a$ and $|o_s|\le B$, so \cref{eq:perturbed-contrast} gives
$r^s_0-r^s_1\ge\frac{\sinh a}{\cosh a+\cosh B}=\tau_0$ for every $s$; apply \cref{lem:gate-floor-completion} with $T_0=T_{\mathrm{gate}}$.
\end{proof}

These results do not show that the acceleration stays negative once the parameters move, or that the floor of
\cref{lem:gate-floor-completion} is reached with $\tau_0\ge1/\mathrm{poly}$. A random-initialization theorem needs one more step: with high
probability, the flow of \cref{lem:factor-flow} reaches such a floor in polynomial time and keeps it for $T_{\mathrm{fin}}$, which requires bounds on
$o_s$, the Gram couplings and $g$ along the trajectory. The head dimension in \cref{prop:gaussian-acceleration} is a proof device: the runs of
\cref{app:planted-perturbation} form the floor at $d_{\mathrm h}=8$.

\paragraph{Matched separation under perturbation.} The perturbed seed of part (i) still satisfies \cref{ass:init}: after division by
$\initscale$, $\|\mathbf W_Q\|,\|\mathbf W_K\|\le\frac98\cdot2\sqrt2$, and embeddings, value map and biases are unchanged. As in \cref{cor:joint-outcome},
\cref{thm:outcome-end} then applies to outcome training from the same perturbed seed with the same trained set, so
\cref{cor:matched-separation} holds with the fixed-perturbation time bound.
